% bab1.tex — dari BAB_I_PENDAHULUAN_POLISHED.md mulai ## A.

\section*{BAB I --- PENDAHULUAN}
\addcontentsline{toc}{section}{BAB I --- PENDAHULUAN}

\subsection*{A. Latar Belakang}

Hadis, selaku sumber kedua ajaran Islam sesudah Al-Qur'an, menduduki posisi pivotal dalam ranah hukum, teologi, maupun etika Islam lewat perannya sebagai \textit{bayan} (klarifikasi), \textit{takhsis} (pembatasan makna umum), serta legislasi otonom atas kasus yang tidak dirinci Al-Qur'an (Musadad, 2026; Razak et al., 2025). Wibawa kenabian itu dikristalkan lewat artikulasi panjang himpunan perdana semisal \textit{al-Muwatta'} beserta \textit{Musannaf}, yang memperlihatkan bagaimana otoritas Nabi dibangun secara metodis oleh ulama generasi pertama (Musadad, 2026; Tanjung \& Fadhlurrahman, 2025). Posisi itu kian menguat dalam fikih masa kini, sebab kritik hadis yang tajam terbukti berkonsekuensi hukum langsung --- umpamanya pada wacana kewarisan Islam, di mana pengujian \textit{matan} dan \textit{sanad} memutus keabsahan dalil (Rozikin, 2026; Bin Jamil, 2024). Pada level epistemologi, hadis menjelma titik temu antara khazanah klasik dan kerangka \textit{maqasid}, sampai-sampai telaahnya melebar dari sekadar verifikasi tekstual menuju orientasi kemaslahatan syariat (Bin Jamil, 2024; Razak et al., 2025). Konsekuensinya, penguasaan hadis secara utuh meniscayakan penelusuran serentak atas dimensi kesejarahan, metodologi, dan kontekstualitasnya, karena setiap klaim normatif yang dinisbatkan kepada Nabi berdampak langsung pada amalan hukum serta moral umat (Rozikin, 2026; Musadad, 2026).

Dari sudut kesejarahan, lintasan hadis terhampar sejak era kenabian hingga zaman modern lewat metamorfosis bertingkat: berawal dari kelisanan, catatan personal (\textit{sahifah}), pembukuan tematik, hingga penetapan kitab induk (Tanjung \& Fadhlurrahman, 2025; Lutfianto et al., 2024). Di masa Nabi, periwayatan berjalan lewat pelbagai wahana --- hafalan, keteladanan praktis, pembelajaran majelis, dan dokumentasi dini selaku penguat ingatan --- sehingga periode tersebut tidak tepat disederhanakan menjadi sepenuhnya oral ataupun sepenuhnya tertulis (Rashwan, 2024; Watukila, 2024). Selepas wafat beliau, hijrah para sahabat ke Syam, Irak, serta Mesir memunculkan transmisi yang bercorak kedaerahan, sementara mobilitas intelektual (\textit{rihlah}) mengarahkan para sahabat menjadi rujukan otoritatif di tiap kawasan (Gil, 2024; Jaiyeoba \& Osmani, 2024). Memasuki generasi tabi'in, jejaring guru-murid antar-kawasan kian solid, simpul-simpul keilmuan regional pun terbentuk, dan \textit{isnad} dimantapkan sebagai kaidah pembuktian yang kian berkesadaran metodis (Watukila, 2024; Lutfianto et al., 2024). Puncaknya pada era ulama klasik, kompilasi semisal \textit{al-Muwatta'}, \textit{al-Musannaf}, dan \textit{Sahihayn} memindahkan tumpuan otoritas dari figur perawi menuju kitab kanon beserta perangkat kritiknya yang telah mapan (Musadad, 2026; Obeid \& Hussein, 2023). Telaah atas penulisan hadis pra-kodifikasi melalui lensa 'Ajjaj al-Khatib menegaskan bahwa peralihan hafalan menuju penulisan berlangsung secara gradual, dengan \textit{qira'ah}, \textit{sama'}, serta audisi berperan selaku jembatan menuju pembukuan menyeluruh (Lutfianto et al., 2024; Tanjung \& Fadhlurrahman, 2025).

Pada rentang lintasan itu mencuat problem akut berupa fabrikasi hadis (\textit{maudu'}) --- kabar buatan yang diklaim bersumber dari Nabi, padahal beliau sama sekali tidak pernah menyabdakannya (Muhamed Ali et al., 2015; Wasman et al., 2023). Batas struktural antara zaman kemurnian periwayatan dan zaman pemalsuan berpangkal pada gejolak politik seusai terbunuhnya Khalifah Usman dan meletusnya \textit{al-fitnah al-kubra}, yang menelurkan kubu politik-teologis yang saling berebut legitimasi (Liew, 2025; Su, 2021). Riset mengenai \textit{afterlife} hadis kenabian seputar khilafah tiga puluh tahun memperlihatkan bahwa satu \textit{matan} yang sama dapat dimaknai-ulang guna menyangga narasi kekuasaan yang bertolak belakang (Liew, 2025; Kara, 2026). Di luar arena politik, pecahnya teologi (Murji'ah, Qadariyah, Jabariyah), rivalitas mazhab fikih, serta ambisi popularitas dan imbalan materiil para tukang cerita (\textit{qussas}) turut menyuburkan industri pemalsuan pada abad kedua--ketiga Hijriah (Wasman et al., 2023; Noor, 2023). Pemetaan jaringan perawi \textit{Sahih Bukhari} menampakkan betapa rumitnya bangunan transmisi itu, sampai-sampai penyimpangan relasional di dalam jejaring \textit{sanad} dapat dibaca sebagai petunjuk adanya distorsi atau sisipan (Alam \& Schneider, 2020; Mosa, 2025). Dengan kata lain, fabrikasi hadis bukan sekadar penyimpangan literatur pinggiran, melainkan krisis teologis berskala besar yang merongrong otentisitas ajaran dari dalam dan mendesak jawaban metodologis yang sistematis (Muhamed Ali et al., 2015; Wasman et al., 2023). Jawaban ulama pun berwujud bangunan penyelamatan yang berpijak pada kerangka \textit{al-jarh wa al-ta'dil}, \textit{takhrij}, serta kritik \textit{matan} (Muhamed Ali et al., 2015; Abbas \& Rawabdeh, 2025). Arti penting \textit{al-jarh wa al-ta'dil} bersumber dari fungsinya selaku penyaring integritas (\textit{'adalah}) serta kekuatan hafalan (\textit{dabt}) perawi, sehingga hanya periwayatan perawi \textit{thiqah} yang diterima sebagai argumen (Muhamed Ali et al., 2015; Su, 2021). Rekonstruksi atas sumbangsih 'Ali b. al-Madini (161--234 H/778--849 M) memperlihatkan bahwa kritik perawi pada generasi dini telah menanam fondasi pembakuan ilmu \textit{'ilal} serta \textit{rijal} yang kelak diteruskan Ibn al-Salah beserta generasi Mamluk (Su, 2021; Noor, 2023). Pada ranah \textit{matan}, kupasan kritis Ibn Kasir terhadap sebagian hadis \textit{Sahihayn} bertema \textit{sirah} di dalam \textit{al-Bidaya wa al-Nihaya} memperlihatkan bahwa otoritas kanonik pun tetap terbuka bagi telaah isi berbasis keselarasan historis-teologis (Calgan, 2024; Ardig et al., 2025). Kajian atas konsep \textit{munkar} pada telaah kritis al-Dzahabi mengenai autentikasi al-Hakim (\textit{al-Mustadrak}) mempertegas bahwa autentikasi bukan putusan tunggal, melainkan dialektika antar-muhaddis yang senantiasa diuji (Abbas \& Rawabdeh, 2025; Obeid \& Hussein, 2023). Lewat penyusunan kitab-kitab \textit{maudu'at}, pembakuan \textit{muqaddimah} ilmu hadis, serta penyaringan sumber yang ketat menurut teladan al-Bukhari, ulama klasik menegakkan daya tahan epistemologis yang melindungi korpus dari penyusupan lanjutan (Obeid \& Hussein, 2023; Noor, 2023).

Memasuki abad kedua puluh satu, hadis berhadapan dengan guncangan global yang dipicu revolusi digital, jejaring sosial, dan kecerdasan artifisial generatif, sehingga transmisi beralih dari \textit{sanad} yang otoritatif menuju otoritas algoritmik dalam ranah siber (Abdulrahman, 2024; Matin Bin Salman, 2025). Hadis kini bukan lagi sekadar teks yang disalurkan, melainkan wacana yang dihasilkan, diedarkan, dan dimaknai-ulang dalam ekosistem digital yang terdesentralisasi (Matin Bin Salman, 2025; Abdulrahman, 2024). Gejala viralitas tanpa validitas menjelma problem gawat manakala cuplikan hadis disebarluaskan secara masif tanpa \textit{takhrij} yang layak, sehingga digitalisasi \textit{takhrij} berubah menjadi keperluan mendesak demi merapatkan laju sirkulasi dengan kecermatan verifikasi (Abdullah et al., 2026; Supriyadi et al., 2020). Riset \textit{living hadith} pada era kecerdasan artifisial menunjukkan bahwa mahasiswa masa kini menjadikan \textit{ChatGPT} sebagai rujukan belajar hadis --- kebiasaan yang merentangkan akses sekaligus menebar bahaya disinformasi keagamaan tanpa bekal literasi kritis yang cukup (Syukur et al., 2026; Supriyadi et al., 2020). Masa depan kajian hadis diwarnai tarik-menarik antara peluang pembakuan korpus digital dan terkikisnya otoritas tradisional; implikasinya, verifikasi otomatis yang dipadukan literasi hadis digital menjadi kebutuhan prioritas (Abdulrahman, 2024; Bin Jamil, 2024).

\textit{State of the art} rentang 2015--2026 memperlihatkan tiga kutub yang saling menggenapi: pemantapan kritik \textit{matan} mutakhir, komputasionalisasi korpus, dan jawaban atas guncangan digital (Razak et al., 2025; Azmi et al., 2019; Abdulrahman, 2024). Pertama, perkembangan kritik teks mencatat migrasi dari hegemoni \textit{sanad} ke arah integrasi \textit{sanad-cum-matan} plus pembacaan kontekstual yang peka \textit{asbab al-wurud} (Razak et al., 2025; Wasman et al., 2023). Telaah perbandingan dampak modernitas terhadap kritik \textit{matan} pada \textit{Fath al-Bari} karya Ibn Hajar serta \textit{Majallat al-Manar} karya Rasyid Rida menampakkan dua kutub, tradisionis dan reformis, dalam menanggapi hadis bertema gejala alam (Al-Imam, 2026; Razak et al., 2025). Kedua, komputasionalisasi melahirkan model \textit{pretrained} bagi pengenal hadis palsu, pemodelan corak kenabian lewat \textit{AraBERT}, dan kerangka graf pengetahuan-\textit{transformer} untuk penguraian ambiguitas narator (Alghamdi et al., 2025; Shaaban et al., 2026; Mosa, 2025). Survei riset berbasis pemrosesan bahasa alami beserta penilaian fitur kebahasaan melalui \textit{machine learning} mengukuhkan capaian akurasi prediksi yang berarti pada pendekatan komputasional, walaupun validasi ulama tetap tak tergantikan (Azmi et al., 2019; Mohamed \& Sarwar, 2022). Ketiga, ramalan otentisitas lewat telaah sentimen, komparasi \textit{deep learning} lawan kompresi lawan pengklasifikasi klasik, dan ontologi \textit{SemanticHadith} memperlihatkan kematangan prasarana semantik verifikasi berskala luas (Haque et al., 2020; Tarmom et al., 2022; Kamran et al., 2023). Kedudukan riset ini ialah mempertemukan telaah historis-filologis klasik dengan agenda komputasional-digital, lewat perangkuman periodisasi, transmisi, dan pemalsuan ke dalam satu narasi yang padu (Tanjung \& Fadhlurrahman, 2025; Abdulrahman, 2024; Mosa, 2025). Berbeda dari riset satu-fase --- misalnya artikulasi wibawa kenabian pada \textit{al-Muwatta'} serta \textit{Musannaf} (Musadad, 2026), atau polemik kesakralan Madinah (Kara, 2026) --- riset ini memetakan lima kurun: Nabi, sahabat, tabi'in, ulama, dan modern, sebagai satu bingkai analitis bagi perkembangan metodologi transmisi beserta kritik (Razak et al., 2025; Lutfianto et al., 2024). Lain halnya dengan riset komputasional yang teknosentris (Alghamdi et al., 2025; Shaaban et al., 2026), riset ini mendudukkan teknologi selaku penerus etos \textit{tathabbut}, \textit{rihlah}, serta \textit{jarh-ta'dil}, bukan substitusi otoritas ulama (Jaiyeoba \& Osmani, 2024; Muhamed Ali et al., 2015; Gil, 2024). Sumbangsihnya mencakup pemadatan historiografis peralihan lisan-tulisan-kodifikasi-digital, taksonomi pemicu pemalsuan berikut bangunan penyelamatannya hingga zaman algoritmik, serta agenda verifikasi digital berlandas paradigma \textit{maqasid} (Bin Jamil, 2024; Abdulrahman, 2024; Kamran et al., 2023). Tulisan ini diharapkan menjadi gerbang yang layak menuju pembahasan periodisasi, penghimpunan, dan pemalsuan hadis pada Bab II, sekaligus mengukuhkan relevansi ilmu hadis selaku disiplin hidup di tengah pergantian zaman (Tanjung \& Fadhlurrahman, 2025; Razak et al., 2025; Abdullah et al., 2026).

\subsection*{B. Rumusan Masalah}

Berdasarkan latar belakang di atas, rumusan masalah penelitian ini adalah:

\begin{enumerate}[leftmargin=1.27cm,itemsep=2pt,parsep=2pt]
\item Bagaimana periodisasi sejarah perkembangan hadis dari masa Nabi hingga era modern?

\item Bagaimana proses penghafalan dan penghimpunan hadis dari tradisi lisan hingga kodifikasi resmi?

\item Apa faktor-faktor yang mendorong pemalsuan hadis dan bagaimana upaya penyelamatannya?

\end{enumerate}
\subsection*{C. Tujuan Penulisan}

Sejalan dengan rumusan masalah, penelitian ini bertujuan:

\begin{enumerate}[leftmargin=1.27cm,itemsep=2pt,parsep=2pt]
\item Mengetahui dan menganalisis periodisasi sejarah perkembangan hadis.

\item Memahami proses penghafalan dan penghimpunan hadis.

\item Menjelaskan faktor-faktor pemalsuan hadis serta upaya penyelamatan yang dilakukan ulama.

\end{enumerate}
\subsection*{D. Manfaat Penulisan}

\begin{itemize}[leftmargin=1.27cm,itemsep=2pt,parsep=2pt]
\item Manfaat teoretis: menambah kekayaan khazanah ilmu hadis melalui perangkuman historiografis atas temuan klasik dan mutakhir, utamanya transisi epistemologis dari kelisanan menuju pembukuan serta digitalisasi (Tanjung \& Fadhlurrahman, 2025; Razak et al., 2025; Abdulrahman, 2024).

\item Manfaat praktis: menyediakan bahan literasi bagi mahasiswa, pengajar, dan publik untuk memeriksa hadis di lingkungan digital, agar terhindar dari \textit{maudu'} sekaligus mampu memakai perangkat verifikasi \textit{takhrij} digital secara kritis (Abdullah et al., 2026; Supriyadi et al., 2020; Syukur et al., 2026).

\end{itemize}
\subsection*{E. Metode Penulisan}

Penelitian ini memakai telaah kepustakaan (\textit{library research}) yang bersifat kualitatif-deskriptif dan historis-analitis. Data meliputi 35 rujukan utama dari jurnal bereputasi: riset periodisasi (Musadad, 2026; Tanjung \& Fadhlurrahman, 2025; Razak et al., 2025), penghimpunan beserta preservasi (Rashwan, 2024; Lutfianto et al., 2024; Jaiyeoba \& Osmani, 2024; Obeid \& Hussein, 2023), kritik berikut pemalsuan (Muhamed Ali et al., 2015; Su, 2021; Abbas \& Rawabdeh, 2025; Mohamed \& Sarwar, 2022), dan kajian digital-kontemporer (Abdulrahman, 2024; Matin Bin Salman, 2025; Alghamdi et al., 2025; Abdullah et al., 2026). Data dihimpun lewat kajian dokumentasi dan perangkuman literatur; analisisnya memakai analisis isi melalui reduksi, pengelompokan tematik merujuk rumusan masalah, dan penarikan simpulan deduktif-analitis.

